% CECS 378 Demo — Exam 1
% Illustrative exam for the lectern worked example.
%
% Questions are written fresh for demo purposes and do NOT reproduce any
% live exam bank.  Topics: symmetric/asymmetric crypto, hashing,
% buffer overflows, access control.
%
% Build:
%   reg-exam-build exam.build.yaml        (pack mode: A/B + roster)
%   reg-exam-build exam1.tex              (legacy single-form mode)

\documentclass[11pt]{article}
\usepackage[margin=0.9in,top=0.7in,bottom=0.85in,headheight=14pt,headsep=14pt]{geometry}
\usepackage[T1]{fontenc}
\usepackage[utf8]{inputenc}

% --- Print-optimized font stack ---
\usepackage{charter}
\usepackage[scaled=0.92]{helvet}
\usepackage[scaled=0.95,varl]{inconsolata}
\renewcommand{\familydefault}{\rmdefault}
\linespread{1.06}
\frenchspacing

\usepackage{enumitem}
\usepackage{needspace}
\usepackage[hyphens]{url}
\usepackage[hidelinks]{hyperref}
\usepackage{listings}
\usepackage{microtype}
\usepackage{fancyhdr}
\usepackage{pifont}
\usepackage{xcolor}
\definecolor{keyred}{HTML}{B71C1C}
\definecolor{keygreen}{HTML}{1B5E20}
\definecolor{keygrey}{HTML}{616161}

% MC choice color: correct (green-bold) vs distractor (grey), gated by \ifanswers
\newcommand{\correctchoice}[1]{\ifanswers\textcolor{keygreen}{\textbf{#1}}\else #1\fi}
\newcommand{\wrongchoice}[1]{\ifanswers\textcolor{keygrey}{#1}\else #1\fi}

% --- Build-time stamp ---
\newcount\stamphour
\newcount\stampminute
\stamphour=\time
\divide\stamphour by 60
\stampminute=\time
\count255=\stamphour
\multiply\count255 by 60
\advance\stampminute by -\count255
\newcommand{\padtwo}[1]{\ifnum#1<10 0\fi\the#1}
\newcommand{\buildtime}{\the\year-\padtwo{\month}-\padtwo{\day}~\padtwo{\stamphour}:\padtwo{\stampminute}}

% --- Answer-key toggle ---
\newif\ifanswers
\answersfalse
\makeatletter
\@ifundefined{answersmode}{}{\answerstrue}
\makeatother

% --- Exam serial ---
\providecommand{\examserial}{XXXXXXXX}

% --- Per-student macros (A.2 doctrine) ---
\makeatletter
\@ifundefined{studentname}{\def\studentname{}}{}
\@ifundefined{studentid}{\def\studentid{}}{}
\@ifundefined{studentserial}{\def\studentserial{}}{}
\makeatother

% \fieldline{width}{value}
\newcommand{\fieldline}[2]{%
  \rlap{\hspace{0.5em}\raisebox{2.2pt}{#2}}\rule[-3pt]{#1}{0.8pt}%
}

% \identityinstruction — pre-filled vs blank
\makeatletter
\newcommand{\identityinstruction}{%
  \ifx\studentname\@empty
    {\sffamily\bfseries\small\ding{43}~~PRINT CLEARLY. UNNAMED EXAMS CANNOT BE RETURNED OR GRADED.}%
  \else
    {\sffamily\bfseries\small\ding{43}~~VERIFY YOUR NAME AND STUDENT ID BELOW ARE CORRECT, THEN ADD TODAY'S DATE.}%
  \fi%
}
\makeatother

% --- Running header ---
\pagestyle{fancy}
\fancyhf{}
\fancyhead[L]{\textit{\small CECS 378 \textperiodcentered\ Computer Security \textperiodcentered\ Su26}}
\fancyhead[R]{\textit{\small Exam 1}}
\fancyfoot[L]{\textit{\footnotesize Generated \buildtime}}
\fancyfoot[C]{\textit{\small p.~\thepage}}
\makeatletter
\fancyfoot[R]{\textit{\footnotesize Serial \texttt{\examserial}%
  \ifx\studentserial\@empty\else
    \ \textperiodcentered\ ID \texttt{\studentserial}%
  \fi}}
\makeatother
\renewcommand{\headrulewidth}{0.4pt}
\renewcommand{\footrulewidth}{0pt}

% --- Letterpress double rule ---
\newcommand{\doublerule}{%
  \noindent\rule{\linewidth}{1.0pt}\par\vspace{1.5pt}%
  \noindent\rule{\linewidth}{0.3pt}\par}

% --- Drafting-style metadata strip ---
\newcommand{\metastrip}[4]{%
  \par\noindent\rule{\linewidth}{0.3pt}\par\vspace{2pt}%
  \noindent\begin{minipage}{\linewidth}%
  \begin{tabbing}%
  \hspace{6em}\=\hspace{16em}\=\hspace{6em}\=\kill%
  \textsc{course}~~\>#1\>\textsc{term}~~\>#2\\[1pt]%
  \textsc{exam}~~\>#3\>\textsc{points}~~\>#4%
  \end{tabbing}%
  \end{minipage}\par\vspace{-4pt}%
  \noindent\rule{\linewidth}{0.3pt}\par}

% --- Section heading ---
\newcommand{\examsection}[1]{%
  \par\vspace{1.4em}%
  \noindent\rule{\linewidth}{0.3pt}\par\vspace{2pt}%
  \noindent{\textsc{\large #1}}%
  \par\vspace{0.4em}}

% --- Writelines helper ---
\newcounter{writelinesi}
\newcommand{\writelines}[1]{%
  \par\vspace{4pt}%
  \setcounter{writelinesi}{0}%
  \loop\ifnum\value{writelinesi}<#1
    \par\noindent\rule{\linewidth}{0.4pt}\vspace{1.4em}%
    \stepcounter{writelinesi}%
  \repeat
  \par\vspace{4pt}}

% --- List metrics ---
\setlist[enumerate,1]{leftmargin=2.4em, itemsep=0.45em, topsep=0.45em, parsep=0pt, label=\textbf{\arabic*.}}
\setlist[enumerate,2]{leftmargin=1.8em, itemsep=0pt, topsep=0.1em, parsep=0pt, label=(\alph*)}

% --- Code block style ---
\lstset{
  basicstyle=\ttfamily\small,
  breaklines=true,
  breakatwhitespace=true,
  columns=fullflexible,
  showstringspaces=false,
  numbers=none,
  tabsize=3,
  aboveskip=4mm,
  belowskip=4mm,
  frame=tb,
  framerule=0.4pt,
  framesep=6pt,
  xleftmargin=12pt,
  xrightmargin=12pt,
}

% --- Even-page padding ---
\AtEndDocument{%
  \clearpage%
  \ifodd\value{page}\else%
    \thispagestyle{plain}%
    \null\vfill%
    \begin{center}\textit{\small This page intentionally left blank.}\end{center}%
    \vfill\null%
  \fi%
}

\setlength{\emergencystretch}{2em}
\sloppy

% Key banner
\ifanswers\fancyhead[R]{\textit{\small Exam 1 --- \textcolor{keyred}{\textbf{KEY}}}}\fi

\date{}

\begin{document}

\metastrip{CECS 378}{Su26}{Exam 1}{50}
\par\vspace{0.4em}
\begin{center}{\Large\itshape Exam 1 --- Introduction to Computer Security}\end{center}
\par\vspace{-0.2em}
\doublerule
\par\vspace{0.9em}

\par\vspace{6pt}
\noindent\identityinstruction
\par\vspace{12pt}

\noindent{\sffamily\bfseries\large NAME:}~\fieldline{13.5cm}{\large\studentname}
\par\vspace{14pt}
\noindent{\sffamily\bfseries\large STUDENT ID\#:}~\fieldline{5.5cm}{\large\ttfamily\studentid}\hfill{\sffamily\bfseries\large DATE:}~\rule[-3pt]{4cm}{0.8pt}
\par\vspace{6pt}

\ifanswers\par\vspace{0.4em}\noindent\colorbox{keyred}{\parbox{\dimexpr\textwidth-2\fboxsep\relax}{\centering\color{white}\sffamily\bfseries\large $\star$~~ANSWER KEY~~---~~NOT FOR DISTRIBUTION~~$\star$}}\par\vspace{0.4em}\fi

\par\vspace{0.6em}
\noindent\ding{43}~\textbf{Directions.}~\textit{Write on the exam paper. You may use one 3''$\times$5'' card of notes (both sides, handwritten only). \textbf{No electronic devices of any kind are allowed on your desk during the exam.}}
\par\vspace{0.4em}

% ============================================================
%  SECTION 1: MULTIPLE CHOICE  (5 questions × 4 pts = 20 pts)
% ============================================================
\examsection{Section 1 \textnormal{\small --- Multiple Choice (4 pts each, 20 pts total)}}
\begin{enumerate}

% --- Q1: Symmetric vs. asymmetric ---
\item \textit{(4 pts)}~Which of the following best describes a \textbf{symmetric-key cipher}?
  \begin{enumerate}[label=(\alph*)]
    \item \wrongchoice{The sender encrypts with a public key; the receiver decrypts with a private key.}
    \item \correctchoice{The same secret key is used for both encryption and decryption.}
    \item \wrongchoice{No key is required; security relies entirely on algorithm secrecy.}
    \item \wrongchoice{A digital signature is produced as a side effect of encryption.}
  \end{enumerate}
  \ifanswers \par\smallskip\textcolor{keyred}{\textbf{Answer: (b).}}
    Symmetric ciphers (e.g., AES) use a single shared secret for both directions.
    (a) describes asymmetric encryption; (c) is ``security through obscurity,'' which is
    not the symmetric model; (d) conflates encryption with signing.\par\smallskip\fi

% --- Q2: Hash property ---
\item \textit{(4 pts)}~A cryptographic hash function is said to be \textbf{collision resistant} if:
  \begin{enumerate}[label=(\alph*)]
    \item \wrongchoice{Given a hash value, it is computationally infeasible to find any input that produces it.}
    \item \wrongchoice{The hash output is exactly the same length as the input.}
    \item \correctchoice{It is computationally infeasible to find two distinct inputs that produce the same hash output.}
    \item \wrongchoice{The hash function runs in constant time regardless of input length.}
  \end{enumerate}
  \ifanswers \par\smallskip\textcolor{keyred}{\textbf{Answer: (c).}}
    Collision resistance means no adversary can efficiently find $x \ne y$ s.t.\ $H(x)=H(y)$.
    (a) is preimage resistance (a different property); (b) is false — hashes compress
    variable-length input to fixed-length output; (d) is unrelated to security.\par\smallskip\fi

% --- Q3: RSA / asymmetric ---
\item \textit{(4 pts)}~In RSA, a message is encrypted for a recipient by using:
  \begin{enumerate}[label=(\alph*)]
    \item \wrongchoice{The sender's private key.}
    \item \wrongchoice{A shared session key negotiated via Diffie-Hellman.}
    \item \wrongchoice{The recipient's private key.}
    \item \correctchoice{The recipient's public key.}
  \end{enumerate}
  \ifanswers \par\smallskip\textcolor{keyred}{\textbf{Answer: (d).}}
    RSA encryption uses the \emph{recipient's public key}; only the recipient's
    matching private key can decrypt. (a) describes signing, not encryption.\par\smallskip\fi

% --- Q4: Buffer overflow ---
\item \textit{(4 pts)}~A classic \textbf{stack buffer overflow} typically overwrites which
  memory region to redirect control flow?
  \begin{enumerate}[label=(\alph*)]
    \item \wrongchoice{The heap allocation header.}
    \item \wrongchoice{The process's environment variables.}
    \item \correctchoice{The saved return address on the stack frame.}
    \item \wrongchoice{The kernel's system-call dispatch table.}
  \end{enumerate}
  \ifanswers \par\smallskip\textcolor{keyred}{\textbf{Answer: (c).}}
    By writing past the end of a stack-allocated buffer, an attacker can overwrite the
    saved return address (and optionally the saved frame pointer) to redirect execution
    to attacker-controlled code or a gadget chain.\par\smallskip\fi

% --- Q5: Access control ---
\item \textit{(4 pts)}~The \textbf{principle of least privilege} states that:
  \begin{enumerate}[label=(\alph*)]
    \item \wrongchoice{A subject should be granted the maximum privileges needed for any possible future task.}
    \item \wrongchoice{All users must share a single privileged account to simplify auditing.}
    \item \wrongchoice{Privileges are determined by a mandatory label assigned at object creation time.}
    \item \correctchoice{A subject should be granted only the minimum privileges necessary to perform its current task.}
  \end{enumerate}
  \ifanswers \par\smallskip\textcolor{keyred}{\textbf{Answer: (d).}}
    Least privilege limits the damage a compromised process or user can cause.
    (a) violates the principle; (b) is the opposite of good practice;
    (c) describes mandatory access control labels, not least privilege.\par\smallskip\fi

\end{enumerate}

% ============================================================
%  SECTION 2: SHORT ANSWER  (3 questions × 10 pts = 30 pts)
% ============================================================
\examsection{Section 2 \textnormal{\small --- Short Answer (10 pts each, 30 pts total)}}
\begin{enumerate}[resume]

% --- Q6: Symmetric vs. asymmetric performance ---
\item \textit{(10 pts)}~Briefly explain \textbf{why symmetric ciphers are generally used to
  encrypt bulk data} in practice, while asymmetric ciphers are typically used only for
  key exchange or digital signatures.  What fundamental trade-off motivates this hybrid
  design?

  \ifanswers
    \par\smallskip\textit{Key answer:}
    Symmetric ciphers (AES, ChaCha20) operate on fixed-size blocks or streams using
    simple arithmetic operations and are therefore orders of magnitude faster than
    asymmetric algorithms.  Asymmetric ciphers (RSA, ECDH) rely on hard mathematical
    problems (integer factorization, discrete logarithm) whose computations are slow —
    encrypting megabytes of data with RSA directly is impractical.  The hybrid design
    uses the asymmetric algorithm only once (to securely exchange a symmetric session
    key), then uses the fast symmetric cipher for all subsequent bulk data.  This gives
    both the key-distribution benefit of public-key cryptography and the performance of
    symmetric encryption.

    \textbf{Partial credit:} award 5--7 pts for identifying the speed difference and
    that hybrid design exists, even without full mechanistic explanation.
    \par\smallskip
  \else
    \writelines{7}
  \fi

% --- Q7: Buffer overflow mitigation ---
\item \textit{(10 pts)}~Name and briefly describe \textbf{two} OS/compiler-level
  mitigations for stack buffer overflow attacks.  For each mitigation, also state
  one way an attacker might attempt to circumvent it.

  \ifanswers
    \par\smallskip\textit{Key answer (any two of the following; 5 pts each):}
    \begin{itemize}[itemsep=2pt]
      \item \textbf{Stack canary.} The compiler inserts a random sentinel value
        (canary) between local buffers and the saved return address.  Before a function
        returns, the canary is checked; a mismatch aborts the process.
        \textit{Circumvention:} leak the canary value via a format-string or
        information-disclosure vulnerability, then overwrite it with the correct value
        when crafting the overflow payload.
      \item \textbf{Non-executable stack (NX / W\^{}X).} Pages are marked either
        writable or executable, never both.  Shellcode injected onto the stack cannot
        run.
        \textit{Circumvention:} return-oriented programming (ROP) — chain together
        existing executable code gadgets (``ret'' sequences in libraries) to build
        arbitrary computations without injecting new code.
      \item \textbf{Address Space Layout Randomization (ASLR).} The OS randomizes the
        base addresses of the stack, heap, and shared libraries at each execution.
        Hardcoded jump targets in the exploit become unreliable.
        \textit{Circumvention:} brute-force the address space (feasible on 32-bit);
        leak an address from a running process via a format-string or memory-disclosure
        bug to de-randomize the layout.
    \end{itemize}
    \par\smallskip
  \else
    \writelines{9}
  \fi

% --- Q8: Discretionary vs. mandatory access control ---
\item \textit{(10 pts)}~Compare \textbf{Discretionary Access Control (DAC)} and
  \textbf{Mandatory Access Control (MAC)}.  For each model: (i) who controls access
  decisions, and (ii) give one concrete example of a system or OS feature that uses it.

  \ifanswers
    \par\smallskip\textit{Key answer:}

    \textbf{DAC.}  The \emph{resource owner} (typically the user who created the
    object) controls access.  The owner sets permissions and may grant them to other
    subjects at will.  Example: Unix/Linux file permissions (\texttt{chmod},
    \texttt{chown}) and POSIX ACLs — a user can \texttt{chmod 777} their own file to
    grant global read/write access.

    \textbf{MAC.}  Access decisions are controlled by the \emph{system policy},
    enforced by the OS kernel (or a reference monitor) and cannot be overridden by
    individual users.  Subjects and objects carry labels (e.g., security clearance
    levels); the policy (Bell-LaPadula, Biba, etc.) determines whether a given
    subject-to-object access is permitted.  Example: SELinux enforces MAC policy
    on RHEL/Fedora — even \texttt{root} cannot read a file if the SELinux policy
    forbids it.

    \textbf{Key distinction:} in DAC a user can accidentally or maliciously share
    sensitive data by relaxing permissions; MAC prevents this because the policy is
    mandatory and users cannot weaken it.

    \textbf{Partial credit:} 3--4 pts per model for a correct description without
    a concrete example.
    \par\smallskip
  \else
    \writelines{9}
  \fi

\end{enumerate}

\par\vfill
\begin{center}\small\textit{End of exam. \quad Total: 50 points.}\end{center}

\end{document}
